% 分类目录：ml
\documentclass[12pt, a4paper]{article}
\usepackage[margin=2.5cm]{geometry}
\usepackage{ctex}
\usepackage{amsmath, amssymb}
\usepackage{listings}
\usepackage{xcolor}
\usepackage{tabularx}
\usepackage{hyperref}

\lstset{
    language=Python,
    basicstyle=\ttfamily\small,
    keywordstyle=\color{blue},
    commentstyle=\color{green!50!black},
    stringstyle=\color{red},
    showstringspaces=false,
    breaklines=true,
    numbers=left,
    numberstyle=\tiny\color{gray},
    frame=single
}

\title{知识点：主成分分析 (PCA)}
\author{}
\date{}

\begin{document}
\maketitle

\section{核心定义}
\textbf{中文定义}：主成分分析（Principal Component Analysis, PCA）是一种最常用的无监督线性降维算法。其通过线性变换将原始数据投影到一组正交的新坐标系中，使得投影后的数据在第一个坐标轴（第一主成分）上的方差最大，第二个坐标轴（第二主成分）上的方差次之，以此类推。

\textbf{英文定义}：Principal Component Analysis (PCA) is an unsupervised linear dimensionality reduction technique that transforms data into a new orthogonal coordinate system such that the greatest variance by some scalar projection of the data comes to lie on the first coordinate (first principal component), the second greatest variance on the second coordinate, and so on.

\textbf{核心直觉}：
找一个投影方向，使得投影后的点"散得最开"，从而尽可能多地保留原始数据的分布信息。

\section{数学推导：最大方差视角}

\subsection{问题建模}
设定原始数据集为 $X = [x_1, x_2, \dots, x_N] \in \mathbb{R}^{d \times N}$，且数据已做\textbf{零均值化}处理（即 $\sum x_i = 0$）。
我们要找一个投影向量 $w$（满足 $\|w\|=1$），使得样本点投射到 $w$ 后的方差最大。
投影后的样本点为 $w^\top x_i$，其方差为：
$$ \text{Var} = \frac{1}{N} \sum_{i=1}^N (w^\top x_i)^2 = w^\top \left( \frac{1}{N} \sum_{i=1}^N x_i x_i^\top \right) w $$
令协方差矩阵 $C = \frac{1}{N} XX^\top$，目标变为最小化：
$$ \max_{w} w^\top C w \quad \text{s.t.} \quad w^\top w = 1 $$

\subsection{拉格朗日乘子法求解}
引入拉格朗日乘子 $\lambda$：
$$ L(w, \lambda) = w^\top C w - \lambda (w^\top w - 1) $$
对 $w$ 求导并令其为 0：
$$ \frac{\partial L}{\partial w} = 2Cw - 2\lambda w = 0 \implies Cw = \lambda w $$
\textbf{结论}：投影方向 $w$ 必须是协方差矩阵 $C$ 的\textbf{特征向量}。代入原式发现方差就是特征值 $\lambda$。因此，最大方差对应的方向就是最大特征值对应的特征向量。

\section{另一种视角：最小重构误差}
PCA 也可以理解为寻找一个超平面对原始空间进行投影，使得投影后的样本点恢复到原始空间时的\textbf{均方误差最小}。
$$ \min \sum_{i=1}^N \|x_i - \hat{x}_i\|^2 $$
推导结果与最大方差视角完全等价。

\section{算法步骤}
\begin{enumerate}
    \item \textbf{去中心化}：对所有样本进行 $x_i \leftarrow x_i - \mu$，其中 $\mu$ 为均值。
    \item \textbf{计算协方差矩阵}：$C = \frac{1}{N} XX^\top$。
    \item \textbf{特征值分解}：计算 $C$ 的特征值和特征向量。
    \item \textbf{选择主成分}：将特征值从大到小排序，选择前 $k$ 个特征值对应的特征向量 $w_1, \dots, w_k$。
    \item \textbf{投影降维}：$X' = W^\top X$。
\end{enumerate}

\section{SVD 实现的优势}
在实际工程中，通常不直接对协方差矩阵进行特征分解，而是对原始矩阵 $X$ 进行\textbf{奇异值分解 (SVD)}：
$$ X = U \Sigma V^\top $$
协方差矩阵 $XX^\top = U \Sigma \Sigma^\top U^\top$。可以看出，特征分解中的特征向量就是 SVD 中的左奇异向量 $U$。
\textbf{优点}：
\begin{enumerate}
    \item SVD 不需要显式计算巨大的协方差矩阵（保存 $d \times d$ 矩阵可能溢出内存）。
    \item SVD 的数值计算比特征分解更稳定（精度更高）。
\end{enumerate}

\section{PCA vs. LDA}
\begin{table}[h]
\centering
\begin{tabularx}{\textwidth}{|l|X|X|}
\hline
\textbf{特性} & \textbf{PCA} & \textbf{LDA} \\
\hline
\textbf{是否监督} & 无监督 & 有监督 (利用标签) \\
\hline
\textbf{核心目标} & 最大化全局方差 (保持信息) & 最大化类间方差/类内方差比 (分类) \\
\hline
\textbf{降维限制} & 无，$1 \sim d$ 均可 & 类别数减一 $M-1$ \\
\hline
\end{tabularx}
\end{table}

\section{关键代码片段}
\begin{lstlisting}
import numpy as np
from sklearn.decomposition import PCA
from sklearn.datasets import load_digits

# 加载手写数字数据集 (64维)
X, _ = load_digits(return_X_y=True)

# 初始化 PCA
pca = PCA(n_components=0.95) # 保留 95% 的方差信息
X_reduced = pca.fit_transform(X)

print(f"原始特征数: {X.shape[1]}")
print(f"降维后特征数: {X_reduced.shape[1]}")
print(f"各个主成分解释的方差比: {pca.explained_variance_ratio_}")
\end{lstlisting}

\section{深度面试题}

\textbf{问：PCA 前为什么一定要做中心化（Normalization/Scaling）？}\\
答：1. \textbf{均值归零}：PCA 推导假设协方差矩阵定义在均值为 0 的基础上，否则第一个主成分可能会指向均值方向而非方差方向；2. \textbf{尺度缩放}：PCA 对特征量级敏感。如果某个特征单位是"公里"，另一个是"厘米"，前者方差会远大，导致 PCA 误认为其携带更多信息。

\textbf{问：如何选择主成分个数 $k$？}\\
答：通常使用"方差累计贡献率"：
$$ \eta = \frac{\sum_{i=1}^k \lambda_i}{\sum_{j=1}^d \lambda_j} \geq t $$
常用的阈值 $t$ 设为 $95\%$ 或 $99\%$。

\section{参考文献}
\begin{enumerate}
    \item Pearson, K. (1901). \textit{On lines and planes of closest fit to systems of points in space}.
    \item Bishop, C. M. (2006). \textit{Pattern Recognition and Machine Learning}, Chapter 12.
    \item 同济大学数学系《线性代数》。
\end{enumerate}

\end{document}
